% =====================================================================
\section{Curvas en el espacio: triedro intrínseco}
\label{sec:frenet}
% =====================================================================

La base adaptada a la \emph{órbita} (no a unas coordenadas) es el triedro de
Frenet: tangente, normal y binormal. En ella la aceleración se separa en una
parte que cambia el módulo de $\vb{v}$ y otra que cambia su dirección.

\subsection{Longitud de arco y triedro de Frenet}

\begin{definicion}[Longitud de arco]
\label{def:arco}
Para una curva regular $\vb{r}(t)$ ($\dot{\vb{r}}\neq\vb{0}$), la
\emph{longitud de arco} desde $t_{0}$ es
\begin{equation}
  s(t)=\int_{t_{0}}^{t}\norm{\dot{\vb{r}}(t')}\dd t'
      =\int_{t_{0}}^{t}\sqrt{\dot x^{2}+\dot y^{2}+\dot z^{2}}\,\dd t',
  \qquad \od{s}{t}=\norm{\vb{v}}=v,\qquad \dd s^{2}=\dd\vb{r}\cdot\dd\vb{r}.
  \label{eq:arco}
\end{equation}
\end{definicion}

\begin{observacion}[La longitud como funcional]
Un \emph{funcional} asigna un número a cada función. La longitud de una curva
lo es: la variable es la curva entera.
\begin{enumerate}
  \item \emph{Gráfica $y=y(x)$ en el plano}:
  \[
    L[y]=\int_{a}^{b}\sqrt{1+y'(x)^{2}}\;\dd x,
    \qquad\text{de la forma } \int_{a}^{b}F(x,y,y')\,\dd x \text{ con } F=\sqrt{1+y'^{2}}.
  \]
  \item \emph{Curva paramétrica}:
  \[
    L[\vb{r}]=\int_{t_{0}}^{t_{1}}\norm{\dot{\vb{r}}(t)}\,\dd t
    =\int_{t_{0}}^{t_{1}}\sqrt{\dot x^{2}+\dot y^{2}+\dot z^{2}}\;\dd t .
  \]
  \item \emph{En coordenadas curvilíneas}, por \eqref{eq:metrica}; por ejemplo, sobre la esfera $r=R$:
  \[
    L[\vb{q}]=\int_{t_{0}}^{t_{1}}\sqrt{h_{1}^{2}(\dot q^{1})^{2}+h_{2}^{2}(\dot q^{2})^{2}+h_{3}^{2}(\dot q^{3})^{2}}\;\dd t,
    \qquad
    L[\theta,\varphi]=R\int_{t_{0}}^{t_{1}}\sqrt{\dot\theta^{2}+\sin^{2}\theta\,\dot\varphi^{2}}\;\dd t .
  \]
\end{enumerate}
\emph{Función frente a funcional}: $s(t)$ en \eqref{eq:arco} es una
\emph{función} de $t$ (la curva está fija); $L[\vb{r}]$ es un \emph{funcional}
(cambia la curva y cambia el número). Además $L$ no depende de la
parametrización: con $t=t(\lambda)$,
$\norm{\dd\vb{r}/\dd t}\,\dd t=\norm{\dd\vb{r}/\dd\lambda}\,\dd\lambda$, así que es una
propiedad de la \emph{órbita}, no de la trayectoria.
\end{observacion}

\begin{definicion}[Triedro de Frenet, curvatura, torsión]
\label{def:frenet}
Parametrizando por $s$ y suponiendo $\dd\utau/\dd s\neq\vb{0}$:
\begin{equation}
  \utau\defeq\od{\vb{r}}{s},\qquad
  \kappa\defeq\norm{\od{\utau}{s}},\qquad
  \un\defeq\frac{1}{\kappa}\od{\utau}{s},\qquad
  \ub\defeq\utau\times\un .
  \label{eq:triedro}
\end{equation}
$\kappa$ es la \emph{curvatura} y $\rho_{c}=1/\kappa$ el \emph{radio de
curvatura}. La \emph{torsión} $\tors$ se define en \eqref{eq:frenet} por
$\dd\ub/\dd s=-\tors\,\un$.
\end{definicion}

\begin{lema}[Derivada de un vector unitario]
\label{lem:unitario}
Si $\norm{\uv{a}(s)}=1$ para todo $s$, entonces $\uv{a}\perp\dd\uv{a}/\dd s$.
\end{lema}

\begin{proof}
$0=\od{}{s}(1)=\od{}{s}(\uv{a}\cdot\uv{a})=2\,\uv{a}\cdot\od{\uv{a}}{s}$.
\end{proof}

En particular $\un\perp\utau$, de modo que $(\utau,\un,\ub)$ es una base
ortonormal dextrógira en cada punto de la curva.

\subsection{Fórmulas de Frenet--Serret}

\begin{teorema}[Frenet--Serret]
\label{teo:frenet}
Existe una función $\tors(s)$, la \emph{torsión}, tal que
\begin{equation}
  \resultado{\od{\utau}{s}=\kappa\,\un,\qquad
  \od{\un}{s}=-\kappa\,\utau+\tors\,\ub,\qquad
  \od{\ub}{s}=-\tors\,\un}
  \label{eq:frenet}
\end{equation}
o, en forma matricial,
\begin{equation}
  \od{}{s}\begin{pmatrix}\utau\\ \un\\ \ub\end{pmatrix}
  =\begin{pmatrix}0&\kappa&0\\ -\kappa&0&\tors\\ 0&-\tors&0\end{pmatrix}
  \begin{pmatrix}\utau\\ \un\\ \ub\end{pmatrix}.
\end{equation}
\end{teorema}

\begin{proof}
\begin{enumerate}
  \item \emph{Primera}: es la definición \eqref{eq:triedro}.
  \item \emph{Tercera}. Por el \cref{lem:unitario}, $\ub'\perp\ub$. Además
  $\ub'=(\utau\times\un)'=\utau'\times\un+\utau\times\un'
  =\kappa\,\un\times\un+\utau\times\un'=\utau\times\un'$, que es
  $\perp\utau$. Ortogonal a $\ub$ y a $\utau$, $\ub'$ es paralelo a $\un$:
  $\ub'=-\tors\,\un$ define $\tors$.
  \item \emph{Segunda}. Se descompone $\un'$ en la base:
  $\utau\cdot\un'=(\utau\cdot\un)'-\utau'\cdot\un=0-\kappa=-\kappa$;
  $\un\cdot\un'=0$ (\cref{lem:unitario});
  $\ub\cdot\un'=(\ub\cdot\un)'-\ub'\cdot\un=0+\tors=\tors$.\qedhere
\end{enumerate}
\end{proof}

\begin{clave}[violeta]{Vector de Darboux}
Las tres fórmulas son $\dd\uv{u}/\dd s=\vb{\omega}\times\uv{u}$ para
$\uv{u}\in\{\utau,\un,\ub\}$, con
$\vb{\omega}=\tors\,\utau+\kappa\,\ub$: al recorrer la curva el triedro gira
con «velocidad angular por unidad de longitud» $\vb{\omega}$; $\kappa$ mide
el giro alrededor de $\ub$ (dentro del plano osculador) y $\tors$ el giro
alrededor de $\utau$ (salida del plano). Curva plana $\iff\tors\equiv0$.
\end{clave}

% ---------------------------------------------------------------------
\begin{figure}[t]
\centering
\tdplotsetmaincoords{65}{110}
\begin{tikzpicture}[tdplot_main_coords,scale=1.0]
  \def\R{1.5}\def\c{0.24}
  \pgfmathsetmacro{\L}{sqrt(\R*\R+\c*\c)}
  % eje
  \draw[aux] (0,0,-0.3) -- (0,0,{\c*rad(1080)+0.4});
  \node[peq,anchor=south] at (0,0,{\c*rad(1080)+0.4}) {eje};
  % hélice (tramo trasero más tenue)
  \foreach \ta/\tb in {110/290,470/650,830/1010}{
    \draw[tenue!60,line width=0.7pt] plot[domain=\ta:\tb,samples=60,smooth]
      ({\R*cos(\x)},{\R*sin(\x)},{\c*rad(\x)});}
  \foreach \ta/\tb in {0/110,290/470,650/830,1010/1080}{
    \draw[texto,line width=1.1pt] plot[domain=\ta:\tb,samples=60,smooth]
      ({\R*cos(\x)},{\R*sin(\x)},{\c*rad(\x)});}
  % triedros
  \foreach \tt/\lab in {240/0,960/1}{
    \pgfmathsetmacro{\px}{\R*cos(\tt)}\pgfmathsetmacro{\py}{\R*sin(\tt)}
    \pgfmathsetmacro{\pz}{\c*rad(\tt)}
    \coordinate (P) at (\px,\py,\pz);
    \draw[vec,cian] (P) -- ++({-\R*sin(\tt)/\L},{\R*cos(\tt)/\L},{\c/\L})
       coordinate (T);
    \draw[vec,rosa] (P) -- ++({-cos(\tt)},{-sin(\tt)},0) coordinate (N);
    \draw[vec,ambar] (P) -- ++({\c*sin(\tt)/\L},{-\c*cos(\tt)/\L},{\R/\L})
       coordinate (B);
    \node[punto] at (P) {};
    \ifnum\lab=1
      \node[etiq,anchor=east,inner sep=2pt] at (T) {$\utau$};
      \node[etiq,anchor=north west,inner sep=1pt] at (N) {$\un$};
      \node[etiq,anchor=south,inner sep=2pt] at (B) {$\ub$};
    \fi}
\end{tikzpicture}\hspace{6mm}%
\begin{tikzpicture}[scale=1.0]
  % plano osculador: circunferencia osculatriz y descomposición de a
  \coordinate (P) at (0,0);
  \coordinate (C) at (0,-2.2);
  \draw[violeta!70,line width=0.5pt,dash pattern=on 3pt off 2pt] (C) circle (2.2);
  \draw[texto,line width=1.1pt] plot[domain=-2.3:2.3,samples=60,smooth]
     (\x,{-\x*\x/4.4});
  \node[punto,fill=violeta,inner sep=1pt] at (C) {};
  \node[peq,anchor=north] at (C) {$C$};
  \draw[aux] (P) -- node[pos=0.88,anchor=west,peq]{$\rho_{c}=1/\kappa$} (C);
  \draw[vec,cian] (P) -- (1.5,0) node[anchor=south,etiq,inner sep=2pt]{$\utau$};
  \draw[vec,rosa] (P) -- (0,-0.85) node[anchor=west,etiq,inner sep=2pt]{$\un$};
  \draw[vec,verde] (P) -- (-1.2,-1.5) node[anchor=north east,etiq,inner sep=1pt]{$\vb{a}$};
  \draw[aux] (-1.2,-1.5) -- (-1.2,0);
  \draw[aux] (-1.2,-1.5) -- (0,-1.5);
  \node[peq,anchor=south,inner sep=3pt] at (-0.6,0) {$a_{T}$};
  \node[peq,anchor=west,inner sep=2pt] at (0,-1.5) {$a_{N}$};
  \node[punto] at (P) {};
  \node[etiq,anchor=south west,inner sep=2pt] at (P) {$P$};
\end{tikzpicture}
\caption{Izquierda: triedro de Frenet $(\utau,\un,\ub)$ sobre una hélice
circular; $\un$ apunta siempre hacia el eje y $\ub$ forma un ángulo fijo con
él. Derecha: en el plano osculador, la circunferencia osculatriz (centro
$C$, radio $\rho_{c}$) y la descomposición $\vb{a}=a_{T}\utau+a_{N}\un$ con
$a_{N}=\kappa v^{2}\ge0$; aquí $a_{T}<0$ (la partícula frena).}
\label{fig:frenet}
\end{figure}
% ---------------------------------------------------------------------

\subsection{Componentes intrínsecas de la aceleración}

\begin{proposicion}[Aceleración tangencial y normal]
\label{prop:intrinseca}
Con $v=\dd s/\dd t$,
\begin{equation}
  \resultado{\vb{v}=v\,\utau,\qquad
  \vb{a}=\dot v\,\utau+\kappa v^{2}\,\un=a_{T}\,\utau+a_{N}\,\un,\qquad
  a_{T}=\dot v,\quad a_{N}=\frac{v^{2}}{\rho_{c}} .}
  \label{eq:intrinseca}
\end{equation}
La aceleración está siempre en el plano osculador ($\vb{a}\cdot\ub=0$).
\end{proposicion}

\begin{proof}
$\vb{v}=\od{\vb{r}}{s}\od{s}{t}=v\,\utau$. Derivando y usando la regla de la
cadena y \eqref{eq:frenet}:
$\vb{a}=\dot v\,\utau+v\,\od{\utau}{s}\od{s}{t}=\dot v\,\utau+\kappa v^{2}\,\un$.
\end{proof}

\begin{observacion}[Significado físico]
$a_{T}$ cambia el \emph{módulo} de la velocidad (acelera o frena);
$a_{N}\ge0$ cambia su \emph{dirección} y apunta al centro de curvatura
(aceleración centrípeta). En un movimiento rectilíneo $a_{N}=0$; en uno
uniforme $a_{T}=0$.
\end{observacion}

\subsection{Cálculo con un parámetro arbitrario}

En la práctica la curva viene dada por $\vb{r}(t)$ y $s(t)$ rara vez admite
inversa explícita. No hace falta.

\begin{proposicion}[Fórmulas con parámetro $t$]
\label{prop:formulas-t}
Si $\vb{v}\times\vb{a}\neq\vb{0}$,
\begin{equation}
  \resultado{
  \begin{gathered}
  \utau=\frac{\vb{v}}{v},\qquad
  \ub=\frac{\vb{v}\times\vb{a}}{\norm{\vb{v}\times\vb{a}}},\qquad
  \un=\ub\times\utau,\\[1mm]
  \kappa=\frac{\norm{\vb{v}\times\vb{a}}}{v^{3}},\qquad
  \tors=\frac{(\vb{v}\times\vb{a})\cdot\dot{\vb{a}}}{\norm{\vb{v}\times\vb{a}}^{2}},\qquad
  a_{T}=\frac{\vb{v}\cdot\vb{a}}{v},\qquad
  a_{N}=\frac{\norm{\vb{v}\times\vb{a}}}{v}.
  \end{gathered}}
  \label{eq:formulas-t}
\end{equation}
\end{proposicion}

\begin{proof}
De \eqref{eq:intrinseca}, $\vb{v}\times\vb{a}=v\utau\times(\dot v\utau+\kappa v^{2}\un)=\kappa v^{3}\,\ub$,
lo que da $\ub$ y $\kappa$; y $\vb{v}\cdot\vb{a}=v\dot v$ da $a_{T}$. Para
$\tors$, derivando $\vb{a}$:
$\dot{\vb{a}}=(\cdots)\utau+(\cdots)\un+\kappa v^{2}\,\dot{\un}$ y
$\dot{\un}\cdot\ub=v\,\tors$ por \eqref{eq:frenet}, luego
$(\vb{v}\times\vb{a})\cdot\dot{\vb{a}}=\kappa v^{3}\cdot\kappa v^{3}\tors=\norm{\vb{v}\times\vb{a}}^{2}\tors$.
\end{proof}

\begin{ejemplo}[Circunferencia de radio $R$]
\label{ej:circ}
Con $\theta=s/R$: $\vb{r}(s)=R\bigl(\cos\tfrac{s}{R},\sin\tfrac{s}{R}\bigr)$,
$\utau=\bigl(-\sin\tfrac{s}{R},\cos\tfrac{s}{R}\bigr)$ y
$\od{\utau}{s}=-\tfrac{1}{R}\bigl(\cos\tfrac{s}{R},\sin\tfrac{s}{R}\bigr)$.
Luego $\kappa=1/R$, $\un=-\vb{r}/R$ apunta al centro y $a_{N}=v^{2}/R$: la
aceleración centrípeta conocida.
\end{ejemplo}

\begin{clave}{Método: triedro intrínseco de $\vb{r}(t)$}
\begin{enumerate}[leftmargin=1.4em]
  \item Derivar: $\vb{v}=\dot{\vb{r}}$, $\vb{a}=\ddot{\vb{r}}$ (y $\dot{\vb{a}}$ si se pide $\tors$).
  \item $v=\norm{\vb{v}}$ y $\utau=\vb{v}/v$.
  \item Calcular $\vb{v}\times\vb{a}$ (determinante) y su módulo.
  \item $\ub=\vb{v}\times\vb{a}/\norm{\vb{v}\times\vb{a}}$, $\un=\ub\times\utau$
        \emph{o bien} $\un=(\vb{a}-a_{T}\utau)/a_{N}$.
  \item $\kappa$, $\tors$, $a_{T}$, $a_{N}$ con \eqref{eq:formulas-t}.
  \item \emph{Comprobar}: $\utau\cdot\un=0$, $a_{T}^{2}+a_{N}^{2}=\norm{\vb{a}}^{2}$, y que
        $\un$ apunta hacia la concavidad.
\end{enumerate}
\end{clave}
